\section{Introduction to Ordinary Differential Equations}

\textit{This lecture corresponds to \S1.1-1.3 in the textbook.}

\subsection{What Are Differential Equations?}

An \bold{differential equation}, which I will often abreviate as a \bold{DE}, is an equation containing derivatives of an unknown function. The two main types of DEs are \bold{ordinary differential equations}, or \bold{ODEs}, which are equations involving functions depending on a single independent variable, and \bold{partial differential equations}, or \bold{PDEs}, which are equations involving functions which depend on more than one independent variable. For example, the equation:
\[ \frac{d^2u}{dt^2} + \frac{du}{dt} + u = f(t) \]
is an ODE that depends on the independent variable $t$. The equation:
\[ -\partial_t^2u + \partial_x^2u + \partial_y^2u + \partial_z^2u = 0 \]
is an example of a PDE that depends on multiple independent variables $x$, $y$, $z$, and $t$. The focus of this course will be on ODEs.


\begin{example}
    Consider the following ODE:
\end{example}
    \[ \frac{dy}{dt} = -y(t) + t \]
    How can we determine what this equation represents? From calculus, we know that $dy/dt$ is the derivative of $y(t)$, and that this derivative represents how fast the function $y$ is changing with respect to $t$. Since we have that $dy/dt = -y(t) + t$, the right hand side shows that the slope will decrease as $y$ increases and will increase as $t$ increases.

    Recall that $\frac{dy}{dt}\big|_p$ represents the slope of $y(t)$ passing through the point $p$. We can thus say that the \bold{direction} of $y(t)$ at point $t = t_0$ is equal to $-y(t_0) + t_0$. With this information, we can plot a \bold{direction field} (or \bold{slope field}) for this differential equation:

    \begin{figure}[H]
        \centering 
        \begin{tikzpicture}[scale=0.54, transform shape]
            \draw[->] (\dmin-0.4, 0) -- (\dmax+0.4, 0) node[right] {$t$};
            \draw[->] (0, \dmin-0.4) -- (0, \dmax+0.4) node[above] (ylabel) {$y$};
            \pgfmathsetmacro{\dminb}{\dmin+0.5}
	         \foreach \x in {\dmin,\dminb,...,\dmax} {
	         	\foreach \y in {\dmin,\dminb,...,\dmax} {
                  \pgfmathsetmacro{\m}{\x-\y}
			         \pgfmathsetmacro{\dt}{(\len/2)/sqrt(1+\m*\m)}
		      	   \pgfmathsetmacro{\dy}{\m*\dt}
			         \draw[gray, very thick] (\x-\dt,\y-\dy) -- (\x+\dt,\y+\dy);
		            }
	         }
         \end{tikzpicture}
         \caption{The slope field for the equation $\frac{dy}{dt} = -y(t) + t$}
     \end{figure}
  
     Remark that, in accordance with our intuition, the slope decreases as we move up the $y$ axis and it increases as we move up the $t$ axis. Each line segment in Figure 1 represents the slope of the equation at one of its solutions. Note that if we travel along the line $y(t) = t$, we see that the slope of the solution is always zero. We say that the line $y(t) = t$ is the \bold{equilibrium solution} of the DE and that it does not change with respect to time (i.e. $\frac{dy}{dt} = 0$):

    
     \begin{figure}[H]
        \centering
        \begin{tikzpicture}[scale=0.6, transform shape]
            \draw[->] (\dmin-0.4, 0) -- (\dmax+0.4, 0) node[right] {$t$};
            \draw[->] (0, \dmin-0.4) -- (0, \dmax+0.4) node[above] (ylabel) {$y$};
            \pgfmathsetmacro{\dminb}{\dmin+0.5}
	         \foreach \x in {\dmin,\dminb,...,\dmax} {
	         	\foreach \y in {\dmin,\dminb,...,\dmax} {
                  \pgfmathsetmacro{\m}{\x-\y}
			         \pgfmathsetmacro{\dt}{(\len/2)/sqrt(1+\m*\m)}
		      	   \pgfmathsetmacro{\dy}{\m*\dt}
			         \draw[gray, very thick] (\x-\dt,\y-\dy) -- (\x+\dt,\y+\dy);
		            }
	         }
              \draw[scale=0.5, domain =-10:10, smooth, variable=\x, red, very thick] plot ({\x}, {\x});      
         \end{tikzpicture}
         \caption{The equilibrium solution (in red) to the equation $\frac{dy}{dt} = -y(t) + t$ is $y(t) = t$.}
     \end{figure}

     \begin{example}
    Consider the differential equation $\frac{dy}{dt} = ay$,
    where $a \in \RR$. Plot direction fields for $a = -1$,$0$, and $1$. What is/are the equilibrium solution(s)?
    \end{example}
   \begin{solution}
   When $a = -1$, we have that $\frac{dy}{dt} = -y$, i.e. the slope of the equation is equal to $-y$. Observe that $\frac{dy}{dt}$ does not change when $t$ changes. The direction field is as follows:
   \begin{figure}[H]
    \centering
    \begin{tikzpicture}[scale=0.6, transform shape]
        \draw[->] (\dmin-0.4, 0) -- (\dmax+0.4, 0) node[right] {$t$};
        \draw[->] (0, \dmin-0.4) -- (0, \dmax+0.4) node[above] (ylabel) {$y$};
        \pgfmathsetmacro{\dminb}{\dmin+0.5}
         \foreach \x in {\dmin,\dminb,...,\dmax} {
             \foreach \y in {\dmin,\dminb,...,\dmax} {
              \pgfmathsetmacro{\m}{-\y}
                 \pgfmathsetmacro{\dt}{(\len/2)/sqrt(1+\m*\m)}
                 \pgfmathsetmacro{\dy}{\m*\dt}
                 \draw[gray, very thick] (\x-\dt,\y-\dy) -- (\x+\dt,\y+\dy);
                }
         }   
     \end{tikzpicture}
     \caption{The direction field for $\frac{dy}{dt} = -y$.}
 \end{figure}
    Observe that the slope of the equation is equal to the value of $-y$. For example, when $y = 1$, we have that $\frac{dy}{dt} = -1$. Furthermore note that $\frac{dy}{dt}$ is independent of $t$ and hence does not change as $t$ changes. The equilibrium solution for this system is the x-axis $y = 0$. 

    The direction field for $a = 1$, i.e. $\frac{dy}{dt} = y$, can be plotted in a very similar way.

    When $a = 0$, we have that $\frac{dy}{dt} = 0$. Hence the slopes in the direction field will equal zero for all $t$ and $y$:

    \begin{figure}[H]
        \centering
        \begin{tikzpicture}[scale=0.6, transform shape]
            \draw[->] (\dmin-0.4, 0) -- (\dmax+0.4, 0) node[right] {$t$};
            \draw[->] (0, \dmin-0.4) -- (0, \dmax+0.4) node[above] (ylabel) {$y$};
            \pgfmathsetmacro{\dminb}{\dmin+0.5}
             \foreach \x in {\dmin,\dminb,...,\dmax} {
                 \foreach \y in {\dmin,\dminb,...,\dmax} {
                  \pgfmathsetmacro{\m}{0}
                     \pgfmathsetmacro{\dt}{(\len/2)/sqrt(1+\m*\m)}
                     \pgfmathsetmacro{\dy}{\m*\dt}
                     \draw[gray, very thick] (\x-\dt,\y-\dy) -- (\x+\dt,\y+\dy);
                    }
             }   
         \end{tikzpicture}
         \caption{The direction field for $\frac{dy}{dt} = 0$. }
     \end{figure}

     Since the slope is always zero, any line passing through the direction field will have slope zero. There are hence infinitely many equilibrium  solutions. \qedhere
   \end{solution}

\subsection{Solutions of Differential Equations}

We ``solve'' differential equations of the form $dy/dt = f(t, y)$ by finding a function $y(t)$ that satisfies it.

\begin{example}
    Solve the differential equation $\frac{dy}{dx} = 2y(t)$.
\end{example}
\begin{solution}
    First, rewrite the equation so that all the terms with $y$ appear on one side of the equation and all the terms with $t$ appear on the other:
    \[\frac{dy}{dt} = 2y \implies \frac{dy}{y} = 2dt\]
    Though it may seem bizarre that we separated $dy/dt$ into two components, we are allowed to do so when solving DEs. We will provide more justified reasoning for doing so in the next lecture.
    
    Now we integrate both sides:
    \begin{align*}
    & \int \frac{dy}{y} = \int 2dt \\
    & \ln|y| = 2t + C
    \end{align*}
    Now solve for $y$:
    \begin{align*}
    & \ln|y| = 2t + C \\
    & |y| = e^{2t + C} \\
    &  y(t) = \pm e^{2t}e^C
    \end{align*}
    Note that $e^C$ is a positive constant, so we can replace it with a new constant $A = e^C$ such that:
    \[ y(t) = Ae^{2t} \]
    Lastly let us check our work by integrating $y(t)$ and ensuring it matches the given equation:
    \begin{align*}
    & y' = 2Ae^{2t} \\
    & y' = 2y(t)  \qedhere
    \end{align*}
\end{solution}

Suppose that we are also given an initial condition, for instance suppose for the previous example that $y(0) = 1$. This initial condition uniquely determines the value of the constant and hence the function $y$. With the above example, if we assume $y(0) = 1$, then we find that:
\begin{align*}
    & y(0) = 1 = Ae^{2(0)} \\
    & 1 = A
\end{align*}

\begin{example}
    Solve the ODE $\frac{dy}{dx} = 3y - 2$, given the initial condition $y(0) = 2$.
\end{example}
\begin{solution}
    First, rearrange the equation so that all the terms with $y$ are on one side and all the terms with $x$ are on the other (remember we are allowed to split $\frac{dy}{dx}$ into components):
    \[ \frac{1}{3y - 2}dy = dx \]
    Now integrate both sides:
    \begin{align*}
        & \int \frac{1}{3y - 2}dy = \int dx \\
        & \frac{1}{3} \ln|3y - 2| + C_1 = x + C_2
    \end{align*}
    If we put the constants together we obtain:
    \begin{align*}
        & \ln|3y - 2| = 3x + C \\
        & |3y - 2| = Ae^{3x} \\
        & 3y - 2 = \pm Ae^{3x} \\
        & 3y =  \pm Ae^{3x} + 2 \\
        & y = \frac{\pm Ae^{3x} + 2 }{3}
    \end{align*}
    Now solving for the initial value:
    \begin{align*}
        & y(0) = 2 = \frac{\pm Ae^{3(0)} + 2 }{3} \\
        & 2 = \frac{A + 2}{3} \\
        & 6 = A + 2 \implies A = 4
    \end{align*}
    So a solution to the initial value problem is $y(t) = \frac{4e^{3t} + 2}{3}$.
\end{solution}

\subsection{Classification of Differential Equations}

The \textbf{order} of a DE is the highest derivative in the equation. For example, the order of $y' = y(t)$ is 1, and the order of $y'' + y + t = 0$ is 2. The general form of an order $n$ DE is:
\[ F(x, u(x), u'(x), u''(x), \dots, u^{(n)}(x)) = 0\]
A solution of the ODE $y^{(n)} = f(x,y, y', \dots, y^{(n-1)})$ on the interval $a < x < b$ is a function $\phi$ such that $\phi', \phi'', \dots, \phi^{(n)}$ exists and satisfies $\phi^{(n)} = f(x, \phi, \phi', \dots, \phi^{(n-1)})$ for all $x \in (a, b)$. Note that $\phi$ and all necessary derivatives must exist on $a < x < b$ and $\phi$ must satisfy the differential equation.

For example, consider $y'' + y = 0$, or $y'' = -y$. Solutions to this DE will be of the form $e^{ix}$ or make use of sines and cosines (for instance, it is easy to check that $y = \sin(x) + \cos(x)$ is a solution). However, naturally the question arises: does a solution always exist, if so it is unique? And how do we find it? We will learn how to answer those questions in the next lecture. 

